\section*{Capítulo \ref{ch:g6:solutions-and-solubility} --- Disoluciones y solubilidad}

\begin{solution}{exo:g6:solutions-and-solubility:1}
\omterm{def:g6:solutions-and-solubility:solute}{Soluto}: el azúcar. \omterm{def:g6:solutions-and-solubility:solute}{Disolvente}: el agua. Sí, el \omterm{def:g6:solutions-and-solubility:solute}{disolvente} es el agua: es
una \omterm{def:g6:solutions-and-solubility:solute}{disolución acuosa}.
\end{solution}

\begin{solution}{exo:g6:solutions-and-solubility:2}
Una disolución que ya no puede \omterm{def:g2:mixing-and-dissolving:dissolve}{disolver} más \omterm{def:g6:solutions-and-solubility:solute}{soluto}. El \omterm{def:g6:solutions-and-solubility:solute}{soluto} que se le
añade se queda sin \omterm{def:g2:mixing-and-dissolving:dissolve}{disolver} en el fondo.
\end{solution}

\begin{solution}{exo:g6:solutions-and-solubility:3}
$250 + 20 = \qty{270}{g}$.
\end{solution}

\begin{solution}{exo:g6:solutions-and-solubility:4}
El etanol: \omterm{def:g6:solutions-and-solubility:miscible}{miscible}. El aceite de cocina: \omterm{def:g6:solutions-and-solubility:miscible}{inmiscible}. El vinagre:
\omterm{def:g6:solutions-and-solubility:miscible}{miscible} (es sobre todo agua).
\end{solution}

\begin{solution}{exo:g6:solutions-and-solubility:5}
La propanona (acetona). GHS02: muy inflamable; GHS07: irrita los ojos,
y sus vapores pueden provocar somnolencia o vértigo.
\end{solution}

\begin{solution}{exo:g6:solutions-and-solubility:6}
$360 \times \frac{500}{1000} = \qty{180}{g}$.
\end{solution}

\begin{solution}{exo:g6:solutions-and-solubility:7}
\qty{200}{g} de agua \omterm{def:g2:mixing-and-dissolving:dissolve}{disuelven} como mucho $360 \times \frac{200}{1000} =
\qty{72}{g}$. No: se quedan en el fondo $80 - 72 = \qty{8}{g}$.
\end{solution}

\begin{solution}{exo:g6:solutions-and-solubility:8}
La disolución pesa $225 + 25 = \qty{250}{g}$; $\frac{25}{250} =
\qty{10}{\percent}$.
\end{solution}

\begin{solution}{exo:g6:solutions-and-solubility:9}
En el primer vaso, donde toda la sal se ha disuelto y la disolución
todavía no está saturada.
\end{solution}

\begin{solution}{exo:g6:solutions-and-solubility:10}
El dióxido de carbono estaba disuelto a presión en la botella cerrada.
Al abrirla baja la presión: el agua puede retener menos gas, y el gas
que sobra sale en forma de burbujas.
\end{solution}

\begin{solution}{exo:g6:solutions-and-solubility:11}
$2000 \div 0.04 = 50\,000$ veces más.
\end{solution}

\begin{solution}{exo:g6:solutions-and-solubility:12}
Un gas se \omterm{def:g2:mixing-and-dissolving:dissolve}{disuelve} menos en agua templada: el estanque caliente contiene
menos oxígeno disuelto, y a los peces les falta oxígeno.
\end{solution}

\begin{solution}{pb:g6:solutions-and-solubility:1}
\textbf{1.} \omterm{def:g6:solutions-and-solubility:solute}{Soluto}: la sal. \omterm{def:g6:solutions-and-solubility:solute}{Disolvente}: el agua.

\textbf{2.} El agua contiene tanta sal como puede: ya no se puede
\omterm{def:g2:mixing-and-dissolving:dissolve}{disolver} más sal en ella.

\textbf{3.} La sal sin \omterm{def:g2:mixing-and-dissolving:dissolve}{disolver} del fondo demuestra que la salmuera está
saturada: contiene la mayor cantidad de sal posible.

\textbf{4.} $360 \times 20 = \qty{7200}{g} = \qty{7.2}{kg}$.

\textbf{5.} $20 + 7.2 = \qty{27.2}{kg}$.

\textbf{6.} $\frac{7.2}{27.2} \approx 0.26$: alrededor del
\qty{26}{\percent}.

\textbf{7.} No: el agua podría \omterm{def:g2:mixing-and-dissolving:dissolve}{disolver} \qty{7.2}{kg}. Todavía podría
\omterm{def:g2:mixing-and-dissolving:dissolve}{disolver} $7.2 - 5 = \qty{2.2}{kg}$.

\textbf{8.} $20 - 5 = \qty{15}{kg}$ de agua (\qty{15}{L}).

\textbf{9.} $360 \times 15 = \qty{5400}{g} = \qty{5.4}{kg}$.

\textbf{10.} Sale de la disolución en forma de cristales sólidos:
cristaliza y cae al fondo.

\textbf{11.} Aparecen cristales blancos nuevos que se amontonan en el
fondo del depósito.

\textbf{12.} Han cristalizado $7.2 - 5.4 = \qty{1.8}{kg}$ de sal.
\end{solution}
